\chapter*{Phụ Lục V. Mẫu Thực Hành Cho Gia Đình}
\addcontentsline{toc}{chapter}{Phụ Lục V. Mẫu Thực Hành Cho Gia Đình}
\markboth{Mẫu Thực Hành Gia Đình}{Mẫu Thực Hành Gia Đình}

\section*{Mẫu 1. Bảng Chia Tiền Tháng}
\begin{truyen}{Bốn Ngăn Tối Thiểu}{Four Minimum Categories}
1. Sinh hoạt thiết yếu: ăn uống, điện nước, đi lại, thuốc men cơ bản.\
\textit{(Essentials: food, utilities, transport, basic medicine.)}

2. Học hành và con cái: học phí, sách vở, khoản phát sinh cho việc học.\
\textit{(Education and children: school fees, books, and study-related expenses.)}

3. Quỹ dự phòng: để riêng, không tiêu chéo.\
\textit{(Emergency fund: kept separate, not mixed into daily spending.)}

4. Tiết kiệm dài hơn: khoản nhỏ nhưng phải đều.\
\textit{(Longer-term saving: small but consistent.)}
\end{truyen}

\section*{Mẫu 2. Cuộc Trò Chuyện Cuối Tháng}
\begin{truyen}{Năm Câu Nên Nói Trong Gia Đình}{Five Useful End-of-Month Questions}
1. Tháng này nhà mình tốn nhất ở đâu?\
\textit{(Where did the household spend the most this month?)}

2. Khoản nào là cần thật, khoản nào có thể bớt?\
\textit{(Which expenses were truly necessary, and which can be reduced?)}

3. Mình đã để dành được bao nhiêu?\
\textit{(How much did we manage to save?)}

4. Có khoản nào khiến cả nhà phải lo lại vào tháng sau?\
\textit{(Which expense may trouble the family again next month?)}

5. Tháng tới chỉ sửa đúng một điều gì?\
\textit{(What is the one thing we will correct next month?)}
\end{truyen}

\section*{Mẫu 3. Mười Quy Ước Nhỏ Trong Nhà}
\begin{truyen}{Nếp Nhà Được Giữ Bằng Những Điều Nhỏ}{Household Order Is Preserved Through Small Rules}
1. Không mua món lớn khi chưa bàn bạc.\
\textit{(Do not buy major items without discussion.)}

2. Không coi tiền dạy thêm là tiền để tiêu hết.\
\textit{(Do not treat extra income as money to be fully spent.)}

3. Không lấy quỹ dự phòng để giải quyết các nhu cầu vui thích.\
\textit{(Do not use the emergency fund for pleasure-based wants.)}

4. Mỗi người phải biết tắt điện, khóa nước, giữ đồ dùng bền.\
\textit{(Everyone should turn off lights, shut water taps, and preserve belongings.)}

5. Ăn vừa đủ, không lấy thừa rồi bỏ.\
\textit{(Take only enough food; do not waste what is served.)}

6. Mỗi tuần dành một buổi kiểm lại chi tiêu.\
\textit{(Reserve one time each week to review spending.)}

7. Khi buồn hoặc mệt, không mua sắm ngay.\
\textit{(Do not shop immediately when sad or tired.)}

8. Con cái được giải thích vì sao phải tiết kiệm.\
\textit{(Children should be told why saving matters.)}

9. Không chạy theo nhà khác vì sĩ diện.\
\textit{(Do not imitate other households out of pride.)}

10. Mỗi tháng phải giữ lại dù chỉ một khoản nhỏ.\
\textit{(Every month, keep at least a small amount.)}
\end{truyen}

\section*{Mẫu 4. Mười Hai Câu Nói Để Tự Giữ Mình}
\begin{truyen}{Nói Với Mình Trước Khi Chi}{What to Say Before Spending}
1. Món này có thật sự cần không?\
\textit{(Do I truly need this?)}

2. Nếu không mua hôm nay, có sao không?\
\textit{(What happens if I do not buy it today?)}

3. Món này đổi bằng bao nhiêu giờ làm việc?\
\textit{(How many working hours does this item cost?)}

4. Tôi đang mua vì nhu cầu hay vì cảm xúc?\
\textit{(Am I buying from need or emotion?)}

5. Tôi có đang mua để đỡ ngại với người khác không?\
\textit{(Am I buying this to avoid embarrassment before others?)}

6. Gia đình tôi cần bình yên hay cần món này hơn?\
\textit{(Does my family need peace more than this item?)}

7. Tôi có thể chờ bảy ngày rồi quyết lại không?\
\textit{(Can I wait seven days and then decide again?)}

8. Tôi đã để dành cho quỹ dự phòng chưa?\
\textit{(Have I already contributed to the emergency fund?)}

9. Nếu con tôi nhìn thấy việc này, tôi có muốn nó học theo không?\
\textit{(If my child saw this, would I want the child to imitate it?)}

10. Tôi có đang chiều cơn buồn của mình không?\
\textit{(Am I indulging my sadness?)}

11. Tôi sẽ thấy nhẹ lòng hơn hay nặng lòng hơn sau khi mua?\
\textit{(Will I feel lighter or heavier after buying this?)}

12. Tôi muốn trở thành kiểu người dùng tiền như thế nào?\
\textit{(What kind of person do I want to become in the way I use money?)}
\end{truyen}

\begin{baihoc}
Một gia đình vững không chỉ nhờ thu nhập. Nó vững nhờ có những nguyên tắc đủ rõ để mọi người cùng giữ.
\textit{(A stable family is not sustained by income alone. It is sustained by principles clear enough for everyone to protect together.)}
\end{baihoc}

\begin{ghinhoanh}
Clear household rules reduce future chaos.

Shared discipline protects shared peace.
\end{ghinhoanh}

\ngancach